\documentclass{article}
\usepackage[T1]{fontenc}
\usepackage{lmodern}
\usepackage{graphicx}
\usepackage{amsmath,amssymb}
\usepackage[a4paper,margin=2cm]{geometry}
\usepackage[absolute,overlay]{textpos}
\setlength{\TPHorizModule}{1mm}
\setlength{\TPVertModule}{1mm}
\usepackage[indonesian]{babel}
\usepackage{hyperref}
\setlength{\emergencystretch}{2em}
% unit: Halaman 18

\begin{document}

% === Halaman 18 (llm) ===

\begin{itemize}
    \item $\text{WSFP}$ dipetakan menjadi $\text{ha*e}$.
    \item Dalam Bahasa Inggris, kata yang mungkin untuk $\text{ha*e}$ hanyalah $\text{ha\textcolor{red}{v}e}$, $\text{ha\textcolor{red}{t}e}$, $\text{ha\textcolor{red}{l}e}$, dan $\text{ha\textcolor{red}{z}e}$.
    \item Dengan mencoba mengganti semua $\text{F}$ di dalam cipherteks dengan $\text{v, t, l,}$ dan $\text{z}$, maka huruf yang cocok adalah $\text{v}$ sehingga $\text{WSFP}$ dipetakan menjadi $\text{have}$.
    \item Dengan mengganti $\text{F}$ menjadi $\text{v}$ pada kriptogram $\text{EPYEPOPDZSZUFPO}$ sehingga menjadi $\text{*e*e*e*tat*ve*}$, maka kata yang cocok untuk ini adalah $\text{representatives}$.
\end{itemize}

\end{document}
